\documentclass{article}
\usepackage[utf8]{inputenc}
\usepackage{geometry}
\geometry{
  a4paper,
  total={170mm,257mm},
  left=20mm,
  top=20mm,
}
\usepackage{graphicx}
\usepackage{titling}
\usepackage{booktabs}
\usepackage{listings}
\usepackage{xcolor}
\usepackage{hyperref}
\usepackage{fancyvrb}
\usepackage{float}
\usepackage{amsmath}

\definecolor{darkblue}{rgb}{0.0, 0.0, 0.5}
\RecustomVerbatimEnvironment{verbatim}{Verbatim}{formatcom=\color{darkblue}}
\setlength{\parindent}{0pt}

\title{\textbf{DEEP LEARNING} \\
{\Large \textbf{Report 4: Neural Network}}}
\date{May 2026}

\usepackage{fancyhdr}
\fancypagestyle{plain}{
  \fancyhf{}
  \fancyfoot[L]{\thedate}
  \fancyfoot[R]{\thepage}
  \fancyhead[L]{University of Science and Technology of Hanoi}
  \fancyhead[R]{Deep Learning}
}

\makeatletter
\def\@maketitle{%
  \newpage
  \vskip 1em
  \begin{center}
  \let \footnote \thanks
    {\LARGE \@title \par}%
    \vskip 1em
  \end{center}
  \par
  \vskip 1em}
\makeatother

\begin{document}

\pagestyle{plain}
\maketitle

\noindent\begin{tabular}{@{}ll}
  Student: & Vu Ha Vy -- 2540112 \\
\end{tabular}

\section*{Implementation}

\textbf{Sigmoid activation function}

Each neuron applies the sigmoid activation function:

\[
\sigma(z) = \frac{1}{1 + e^{-z}}
\]

where \(z\) is the linear combination between the inputs and weights.

\begin{verbatim}
def sigmoid(z):
    return 1.0 / (1.0 + math.exp(-z))
\end{verbatim}

\textbf{Neuron implementation}

A neuron stores its weights and bias. During feedforward, the neuron
computes the weighted sum:

\[
z = \sum_{i=1}^{n} x_i w_i + b
\]

and then applies the sigmoid activation:

\[
a = \sigma(z)
\]

\begin{verbatim}
class Neuron:
    def __init__(self, weights, bias):
        self.weights = weights
        self.bias = bias
    
    def feedforward(self, inputs):  
        z = 0
        for i in range(len(inputs)):
            z += inputs[i] * self.weights[i]
        z += self.bias

        a = sigmoid(z)
        return a
\end{verbatim}

\textbf{Layer implementation}

A layer is composed of several neurons. The outputs of all neurons are
stored in a list and returned as the output of the layer.

\begin{verbatim}
class Layer:                  
    def __init__(self, neurons):
        self.neurons = neurons
    
    def feedforward(self, inputs): 
        outputs = []

        for neuron in self.neurons:
            output = neuron.feedforward(inputs)
            outputs.append(output)

        return outputs
\end{verbatim}

\textbf{Neural network implementation}

The neural network is implemented as a sequence of layers.
The output of one layer becomes the input of the next layer. For each layer:

\[
z^{(l)} = (W^{(l)})^T a^{(l-1)} + b^{(l)}
\]

\[
a^{(l)} = \sigma(z^{(l)})
\]

\begin{verbatim}
class NeuralNetwork:            
    def __init__(self, layers):
        self.layers = layers
    
    def feedforward(self, inputs):
        outputs = inputs
        
        for layer in self.layers:
            outputs = layer.feedforward(outputs)
        
        return outputs
\end{verbatim}

\textbf{Neural Network Architecture}

The architecture of the neural network is read from the file
\texttt{nn.txt}. The first line contains the number of layers.
The following lines contain the number of neurons in each layer. For this experiment:

\begin{itemize}
    \item Input layer: 2 neurons
    \item Hidden layer: 2 neurons
    \item Output layer: 1 neuron
\end{itemize}

% \begin{verbatim}
% def read_txt(file_path):
%     with open(file_path, "r") as f:
%         lines = f.readlines()

%     num_layers = int(lines[0].strip())
%     layer_sizes = [int(line.strip()) 
%                    for line in lines[1:num_layers+1]]

%     return num_layers, layer_sizes

% num_layers, layer_sizes = read_txt("nn.txt")
% print(num_layers, layer_sizes)
% \end{verbatim}

The network is initialized using object oriented programming.
Each layer contains neurons, and each neuron contains weights and bias.
The weights and bias are initialized randomly in the interval \([0,1]\).

\begin{verbatim}
def build_nn(layer_sizes):
    layers = []

    for i in range(1, len(layer_sizes)):
        n_inputs  = layer_sizes[i - 1]
        n_neurons = layer_sizes[i]

        neurons = [
            Neuron(
                weights=[random.uniform(0, 1)
                         for _ in range(n_inputs)],
                bias=random.uniform(0, 1)
            )
            for _ in range(n_neurons)
        ]

        layers.append(Layer(neurons))

    return NeuralNetwork(layers)

neural_network = build_nn(layer_sizes)
\end{verbatim}

\section*{XOR Experiment}

The XOR problem cannot be solved using a single logistic regression
classifier because the data is not linearly separable. The XOR truth table is:

\begin{table}[H]
\centering
\begin{tabular}{ccc}
\toprule
$x_1$ & $x_2$ & $x_1 \oplus x_2$ \\
\midrule
0 & 0 & 0 \\
0 & 1 & 1 \\
1 & 0 & 1 \\
1 & 1 & 0 \\
\bottomrule
\end{tabular}
\caption{Truth table of XOR}
\end{table}

The hidden layer is designed using two intermediate logical operations:

\begin{itemize}
    \item NOT($x_1$ AND $x_2$)
    \item $x_1$ OR $x_2$
\end{itemize}

The output layer computes:

\[
\text{XOR} = \text{NOT}(x_1 \land x_2) \land (x_1 \lor x_2)
\]

The weights are initialized exactly as in the lecture slides.

\begin{verbatim}
neural_network.layers[0].neurons[0].weights = [-1, -1]
neural_network.layers[0].neurons[0].bias = 1.5

neural_network.layers[0].neurons[1].weights = [1, 1]
neural_network.layers[0].neurons[1].bias = -0.5

neural_network.layers[1].neurons[0].weights = [1, 1]
neural_network.layers[1].neurons[0].bias = -1.5
\end{verbatim}

\section*{Results}

The network is tested with the four possible XOR inputs.

\begin{verbatim}
data = [[0,0], [0,1], [1,0], [1,1]]

for i in data:
    output = neural_network.feedforward(i)
    print(i, ":", round(output[0], 5))
\end{verbatim}

The obtained outputs are:

\begin{table}[H]
\centering
\begin{tabular}{cccc}
\toprule
Input & Expected Output & Network Output & Predicted Class \\
\midrule
(0, 0) & 0 & 0.42436 & 0 \\
(0, 1) & 1 & 0.43657 & 0 \\
(1, 0) & 1 & 0.43657 & 0 \\
(1, 1) & 0 & 0.42436 & 0 \\
\bottomrule
\end{tabular}
\caption{Results of the XOR neural network}
\end{table}

Using the decision threshold $0.5$, all four outputs are classified as 0.
Therefore, the network gives the correct result for $(0,0)$ and $(1,1)$,
but gives the wrong result for $(0,1)$ and $(1,0)$. This is because the sigmoid activation function does not return exact boolean values 0 and 1. 
For example, for input $(0,1)$, the two hidden neurons return approximately
$0.62246$ and $0.62246$. Although $(0,1)$ should belong to class 1, the final output is still below $0.5$, so it is classified as 0. The output neuron is computed by:

\[
z = 0.62246 + 0.62246 - 1.5 = -0.25508
\]

\[
\hat{y} = \sigma(-0.25508) = 0.43657
\]



\end{document}